\subsection{Roots of Unity}

\frame{

\frametitle{\secname: \subsecname}

\begin{columns}
\column{0.5\textwidth}
\begin{itemize}
\item<1-> The even/odd trick works because repeatedly multiplying by $-1$
\\cycles between $-1$ and $1$
\item<1-> It’d be convenient if we had other values making this kind of cycle, and if these cycles were longer
\item<2-> Complex numbers give us values $r$ where $r^n=1$
\item<2-> These are called roots of unity
\item<2-> \fft{} evaluates $f$ at $1,r,r^2,\dots,r^{n-1}$
\end{itemize}

\column{0.5\textwidth}
\begin{figure}
\centering
\begin{tikzpicture}[scale=1.5,line width=1pt,mark=*,mark size=2pt]
\draw[<->] (-1.5,0) -- (1.5,0);
\only<1>{
\draw (0,-0.25) -- (0,0.25);
}
\only<2->{
\draw[<->] (0,-1.5) -- (0,1.5);
}

\only<1->{
\draw[-{Latex[length=8pt]}] plot[domain=0:180,variable=\t,mark indices=1] ({cos(\t)},{sin(\t)}) node[above left] {$-1$};
\draw[-{Latex[length=8pt]}] plot[domain=180:360,variable=\t,mark indices=1] ({cos(\t)},{sin(\t)}) node [above left] {$1$};
}

\only<2->{
\draw[-{Latex[length=8pt]}] plot[domain=0:90,variable=\t,mark indices=1] ({cos(\t)},{sin(\t)}) node[above left] {$i$};
\draw[-{Latex[length=8pt]}] plot[domain=90:180,variable=\t,mark indices=1] ({cos(\t)},{sin(\t)}) node[above left] {$-1$};
\draw[-{Latex[length=8pt]}] plot[domain=180:270,variable=\t,mark indices=1] ({cos(\t)},{sin(\t)}) node [above left] {$-i$};
\draw[-{Latex[length=8pt]}] plot[domain=270:360,variable=\t,mark indices=1] ({cos(\t)},{sin(\t)}) node [above left] {$1$};
}
\end{tikzpicture}
\end{figure}

\end{columns}

}

